\documentclass{beamer}
\usepackage{tikz,amsmath,hyperref,graphicx,stackrel,animate,media9}
\usetikzlibrary{positioning,shadows,arrows,shapes,calc}
\newcommand{\argmax}{\operatornamewithlimits{argmax}}
\newcommand{\argmin}{\operatornamewithlimits{argmin}}
\mode<presentation>{\usetheme{Frankfurt}}
\AtBeginSection[]
{
  \begin{frame}<beamer>
    \frametitle{Outline}
    \tableofcontents[currentsection,currentsubsection]
  \end{frame}
}
\title{Lecture 16: Discrete-Time Fourier Transform (DTFT)}
\author{Mark Hasegawa-Johnson}
\date{ECE 401: Signal and Image Analysis}  
\begin{document}

% Title
\begin{frame}
  \maketitle
\end{frame}

% Title
\begin{frame}
  \tableofcontents
\end{frame}

%%%%%%%%%%%%%%%%%%%%%%%%%%%%%%%%%%%%%%%%%%%%
\section[Review]{Review: Frequency Response}
\setcounter{subsection}{1}

\begin{frame}
  \frametitle{Response of LSI System to Periodic Inputs}
  Suppose we compute $y[n]=x[n]\ast h[n]$, where
  \begin{align*}
  x[n] &= \frac{1}{N}\sum_{k=0}^{N-1} X[k] e^{j2\pi kn/N},~\mbox{and}\\
  y[n] &= \frac{1}{N}\sum_{k=0}^{N-1} Y[k] e^{j2\pi kn/N}.
  \end{align*}
  The relationship between $Y[k]$ and $X[k]$ is given by the frequency
  response:
  \[
  Y[k] = H(k\omega_0) X[k]
  \]
  where
  \[
  H(\omega) = \sum_{n=-\infty}^\infty h[n]e^{-j\omega n}
  \]
\end{frame}

\begin{frame}
  \frametitle{Response of LSI System to Aperiodic Inputs}

  But what about signals that never repeat themselves?

  Can we still write something like
  \[
  Y(\omega)=H(\omega)X(\omega)?
  \]

\end{frame}
  
%%%%%%%%%%%%%%%%%%%%%%%%%%%%%%%%%%%%%%%%%%%%
\section[DTFT]{Discrete Time Fourier Transform}
\setcounter{subsection}{1}

\begin{frame}
  \frametitle{Aperiodic}
  
  An ``aperiodic signal'' is a signal that is not periodic.
  \begin{itemize}
  \item Music: strings, woodwinds, and brass are periodic, drums and rain sticks are aperiodic.
  \item Speech: vowels and nasals are periodic, plosives and fricatives are aperiodic.
  \item Images: stripes are periodic, clouds are aperiodic.
  \item Bioelectricity: heartbeat is periodic, muscle contractions are aperiodic.
  \end{itemize}
\end{frame}

\begin{frame}
  \frametitle{Periodic}

  The spectrum of a periodic signal is given by its Fourier series.
  In discrete time, that's:
  \begin{align*}
    X_k &= \frac{1}{N_0}\sum_{n=-\frac{N_0}{2}}^{\frac{N_0-1}{2}} x[n] e^{-j\frac{2\pi kn}{N_0}}\\
    &= \frac{1}{N_0}\sum_{n=-\frac{N_0}{2}}^{\frac{N_0-1}{2}} x[n] e^{-j\omega n}
  \end{align*}
  and that gives the frequency content of the signal, at the frequency
  $\omega=\frac{2\pi k}{N_0}$.

  Here I'm using  $n\in\left\{-\frac{N_0}{2},\ldots,\frac{N_0-1}{2}\right\}$,
  but the sum could be over any sequence of $N_0$ continuous samples.
\end{frame}

\begin{frame}
  \frametitle{Aperiodic}

  An aperiodic signal is one that {\bf never} repeats itself.  So we
  want something like the limit, as $N_0\rightarrow\infty$, of the
  Fourier series.  Here is the simplest such thing that is useful:
  \begin{block}{Discrete-Time Fourier Transform (DTFT)}
  \[
  X(\omega) = \sum_{n=-\infty}^\infty x[n]e^{-j\omega n}
  \]
  \end{block}
\end{frame}
    
\begin{frame}
  \frametitle{Fourier Series vs. Fourier Transform}

  The Fourier Series coefficients are:
  \begin{align*}
    X_k &= \frac{1}{N_0}\sum_{n=-\frac{N_0}{2}}^{\frac{N_0-1}{2}} x[n] e^{-j\omega n}
  \end{align*}
  The Fourier transform is:
  \[
  X(\omega) = \sum_{n=-\infty}^\infty x[n]e^{-j\omega n}
  \]
  Notice that, besides taking the limit as $N_0\rightarrow\infty$, we
  also got rid of the $\frac{1}{N_0}$ factor.  So we can think of the
  DTFT as
  \[
  X(\omega) = \lim_{N_0\rightarrow\infty,\omega=\frac{2\pi k}{N_0}} N_0 X_k
  \]
  where the limit is: as $N_0\rightarrow\infty$, and
  $k\rightarrow\infty$, but $\omega=\frac{2\pi k}{N_0}$ remains
  constant.
\end{frame}

\begin{frame}
  \frametitle{Inverse DTFT}

  In order to convert $X(\omega)$ back to $x[n]$, we'll take advantage
  of orthogonality:
  \[
  \int_{-\pi}^\pi e^{j\omega(m-n)} d\omega =
  \begin{cases}
    2\pi & m=n\\
    0 & (m-n)=\mbox{any nonzero integer}
  \end{cases}
  \]
\end{frame}
  
\begin{frame}
  \frametitle{Inverse DTFT}

  Taking advantage of orthogonality, we can see that
  \begin{align*}
    \frac{1}{2\pi}\int_{-\pi}^\pi X(\omega)e^{j\omega m}d\omega\\
    &=\frac{1}{2\pi}\int_{-\pi}^{\pi}
    \left(\sum_{n=-\infty}^\infty x[n]e^{-j\omega n}\right)e^{j\omega m}d\omega\\
    &=\frac{1}{2\pi}\sum_{n=-\infty}^\infty x[n]\int_{-\pi}^\pi e^{j\omega(m-n)}d\omega\\
    = x[m]
  \end{align*}
\end{frame}

  
\begin{frame}
  \frametitle{Fourier Series and Fourier Transform}
  Discrete-Time Fourier Series (DTFS):
  \begin{align*}
    X_k &= \frac{1}{N_0}\sum_{n=0}^{N_0-1} x[n] e^{-j\frac{2\pi kn}{N_0}}\\
    x[n] &= \sum_{k=0}^{N_0-1} X_ke^{j\frac{2\pi kn}{N_0}}
  \end{align*}
  Discrete-Time Fourier Transform (DTFT):
  \begin{align*}
    X(\omega) &= \sum_{n=-\infty}^{\infty} x[n] e^{-j\omega n}\\
    x[n] &= \frac{1}{2\pi}\int_{-\pi}^\pi X(\omega)e^{j\omega n}d\omega
  \end{align*}
  
\end{frame}

\begin{frame}
  \frametitle{Quiz}

  Go to the course webpage, try the quiz!
\end{frame}

%%%%%%%%%%%%%%%%%%%%%%%%%%%%%%%%%%%%%%%%%%%%%%%%%%%%%%%%%%%%%%%%%%%%%%%%%%%%%%%%%%%%
\section[DTFT Properties]{Properties of the DTFT}
\setcounter{subsection}{1}

\begin{frame}
  \frametitle{Properties of the DTFT}

  In order to better understand the DTFT, let's discuss these properties:
  \begin{enumerate}
    \setcounter{enumi}{-1}
  \item Periodicity
  \item Linearity
  \item Time Shift
  \item Frequency Shift
  \item Filtering is Convolution
  \end{enumerate}
  Property \#4 is actually the reason why we invented the DTFT in the first place.
  Before we discuss it, though, let's talk about the others.
\end{frame}

\begin{frame}
  \frametitle{0. Periodicity}

  The DTFT is periodic with a  period of $2\pi$.  That's just because  $e^{j2\pi}=1$:
  \begin{align*}
    X(\omega) &= \sum_n x[n]e^{-j\omega n}\\
    X(\omega+2\pi) &= \sum_n x[n]e^{-j(\omega+2\pi) n} = \sum_n x[n]e^{-j\omega n} = X(\omega)\\
    X(\omega-2\pi) &= \sum_n x[n]e^{-j(\omega-2\pi) n} = \sum_n x[n]e^{-j\omega n} = X(\omega)
  \end{align*}
  For example, the inverse DTFT can be defined in two different ways:
  \[
  x[n]=\frac{1}{2\pi}\int_{-\pi}^\pi X(\omega)e^{j\omega n}d\omega =
  \frac{1}{2\pi}\int_{0}^{2\pi} X(\omega)e^{j\omega n}d\omega
  \]
  Those two integrals are equal because $X(\omega+2\pi)=X(\omega)$.
\end{frame}

\begin{frame}
  \frametitle{1. Linearity}

  The DTFT is linear:
  \[
  z[n] = ax[n]+by[n]~~~\leftrightarrow~~~
  Z(\omega)=aX(\omega)+bY(\omega)
  \]
  {\bf Proof:}
  \begin{align*}
    Z(\omega) &= \sum_n z[n]e^{-j\omega n}\\
    &= a\sum_n x[n]e^{-j\omega n} + b\sum_n y[n]e^{-j\omega n}\\
    &= aX(\omega) + bY(\omega)
  \end{align*}
\end{frame}

\begin{frame}
  \frametitle{2. Time Shift Property}

  Shifting in time is the same as multiplying by a  complex exponential in frequency:
  \[
  z[n] = x[n-n_0]~~~\leftrightarrow~~~
  Z(\omega)=e^{-j\omega n_0}X(\omega)
  \]
  {\bf Proof:}
  \begin{align*}
    Z(\omega) &= \sum_{n=-\infty}^{\infty} x[n-n_0]e^{-j\omega n}\\
    &= \sum_{m=-\infty}^{\infty} x[m]e^{-j\omega (m+n_0)}~~~\left(\mbox{where}~m=n-n_0\right)\\
    &= e^{-j\omega n_0} X(\omega)
  \end{align*}
\end{frame}

\begin{frame}
  \frametitle{3. Frequency Shift Property}

  Shifting in frequency is the same as multiplying by a complex exponential in time:
  \[
  z[n] = x[n]e^{j\omega_0 n}~~~\leftrightarrow~~~
  Z(\omega)=X(\omega-\omega_0)
  \]
  {\bf Proof:}
  \begin{align*}
    Z(\omega) &= \sum_{n=-\infty}^{\infty} x[n]e^{j\omega_0 n}e^{-j\omega n}\\
    &= \sum_{n=-\infty}^{\infty} x[n]e^{-j(\omega-\omega_0) n}\\
    &= X(\omega-\omega_0)
  \end{align*}
\end{frame}

\begin{frame}
  \frametitle{4. Convolution Property}

  Convolving in time is the same as multiplying in frequency:
  \[
  y[n]=h[n]\ast x[n]~~~\leftrightarrow
  Y(\omega)=H(\omega)X(\omega)
  \]
  {\bf Proof:}
  Remember that $y[n]=h[n]\ast x[n]$ means that $y[n]=\sum_{m=-\infty}^\infty h[m]x[n-m]$.  Therefore,
  \begin{align*}
    Y(\omega) &= \sum_{n=-\infty}^{\infty}\left(\sum_{m=-\infty}^\infty h[m]x[n-m]\right)e^{-j\omega n}\\
    &= \sum_{m=-\infty}^{\infty}\sum_{n=-\infty}^\infty\left(h[m]x[n-m]\right)e^{-j\omega m}e^{-j\omega (n-m)}\\
    &= \left(\sum_{m=-\infty}^{\infty}h[m]e^{-j\omega m}\right)
    \left(\sum_{(n-m)=-\infty}^\infty  x[n-m]e^{-j\omega (n-m)}\right)\\
    &= H(\omega)X(\omega)
  \end{align*}
\end{frame}

%%%%%%%%%%%%%%%%%%%%%%%%%%%%%%%%%%%%%%%%%%%%
\section[Averaging]{The Local Averaging Filters}
\setcounter{subsection}{1}

\begin{frame}
  \frametitle{Local Average Filters}

  Let's go back to the local averaging filter. I want to define two 
  different types of local average: centered, and delayed.
  \begin{itemize}
  \item {\bf Centered local average:} This one averages
    $\left(\frac{L-1}{2}\right)$ future samples,
    $\left(\frac{L-1}{2}\right)$ past samples, and $x[n]$:
    \[
    y_c[n] = \frac{1}{L}\sum_{m=-\left(\frac{L-1}{2}\right)}^{\left(\frac{L-1}{2}\right)} x[n-m]
    \]
  \item {\bf Causal local average:} This one averages $x[n]$ and $L-1$ of its
    past samples:
    \[
    y_d[n] = \frac{1}{L}\sum_{m=0}^{L-1} x[n-m]
    \]
  \end{itemize}
  Notice that $y_d[n] = y_c\left[n-\left(\frac{L-1}{2}\right)\right]$.
\end{frame}

\begin{frame}
  \frametitle{Local Average Filters}

  We can write both of these as filters:
  \begin{itemize}
  \item {\bf Centered local average:}
    \begin{align*}
      y_c[n] &= f_c[n]\ast x[n]\\
      f_c[n] &= \begin{cases} \frac{1}{L}& -\left(\frac{L-1}{2}\right)\le n\le\left(\frac{L-1}{2}\right)\\
        0&\mbox{otherwise}\end{cases}
    \end{align*}
  \item {\bf Causal local average:}
    \begin{align*}
      y_d[n] &= f_d[n]\ast x[n]\\
      f_d[n] &= \begin{cases} \frac{1}{L}& 0\le n\le L-1\\
        0&\mbox{otherwise}\end{cases}
    \end{align*}
  \end{itemize}
\end{frame}

\begin{frame}
  \frametitle{Local Average Filters}
  \centerline{\includegraphics[height=2.5in]{exp/localaveragefilters.png}}
  Notice that $f_d[n]=f_c\left[n-\left(\frac{L-1}{2}\right)\right]$.
\end{frame}  

\begin{frame}
  \frametitle{The relationship between centered local average and delayed local average}

  Notice that $f_d[n]=f_c[n-\frac{L-1}{2}]$.  We can find the
  relationship between their DTFTs using variable substitution, with
  the variable $m=n-\frac{L-1}{2}$:
  \begin{align*}
    F_d(\omega) &= \sum_{n=-\infty}^\infty f_d[n]e^{-j\omega n}\\
    &=\sum_{n=-\infty}^\infty f_c[n-\frac{L-1}{2}]e^{-j\omega n}\\
    &=e^{-j\omega\frac{L-1}{2}}\sum_{m=-\infty}^\infty f_c[m]e^{-j\omega m}\\
    &=e^{-j\omega\frac{L-1}{2}}F_c(\omega)
  \end{align*}
\end{frame}
  
\begin{frame}
  \frametitle{The frequency response of a local average filter}

  Let's find the frequency response of
  \[
  f_d[n] = \begin{cases} \frac{1}{L}& 0\le m\le L-1\\
    0&\mbox{otherwise}\end{cases}
  \]
  The formula is
  \[
  F_d(\omega) = \sum_m f[m]e^{-j\omega m},
  \]
  so,
  \[
  F_d(\omega) = \sum_{m=0}^{L-1} \frac{1}{L}e^{-j\omega m}
  \]
\end{frame}
  
\begin{frame}
  \frametitle{The frequency response of a local average filter}

  \[
  F_d(\omega) = \sum_{m=0}^{L-1} \frac{1}{L}e^{-j\omega m}
  \]
  This is just a standard geometric series,
  \[
  \sum_{m=0}^{L-1} a^m = \frac{1-a^L}{1-a},
  \]
  so:
  \[
  F_d(\omega) = \frac{1}{L}\left(\frac{1-e^{-j\omega L}}{1-e^{-j\omega}}\right)
  \]
\end{frame}

\begin{frame}
  \frametitle{The frequency response of a local average filter}

  We now have an extremely  useful transform pair:
  \[
  f_d[n] = \begin{cases} \frac{1}{L}& 0\le m\le L-1\\
    0&\mbox{otherwise}\end{cases}~~~\leftrightarrow~~~
  F_d(\omega) = \frac{1}{L}\left(\frac{1-e^{-j\omega L}}{1-e^{-j\omega}}\right)
  \]
  Let's attempt to convert that into polar form, so we can find
  magnitude and phase response. Notice that both the numerator and the
  denominator are subtractions of complex numbers, so we might be able
  to use $2j\sin(x)=e^{jx}-e^{-jx}$ for some $x$.  Let's try:
  \begin{align*}
    \frac{1}{L}\left(\frac{1-e^{-j\omega L}}{1-e^{-j\omega}}\right)
    &= \frac{1}{L}\frac{e^{-j\omega L/2}}{e^{-j\omega/2}}
    \left(\frac{e^{j\omega L/2}-e^{-j\omega L/2}}{e^{j\omega/2}-e^{-j\omega/2}}\right)\\
    &= e^{-j\omega\left(\frac{L-1}{2}\right)}\frac{1}{L}
    \left(\frac{2j\sin(\omega L/2)}{2j\sin(\omega/2)}\right)\\
    &= e^{-j\omega\left(\frac{L-1}{2}\right)}\frac{1}{L}
    \left(\frac{\sin(\omega L/2)}{\sin(\omega/2)}\right)
  \end{align*}
\end{frame}
  
\begin{frame}
  \frametitle{The frequency response of a local average filter}

  Now we have $F_d(\omega)$ in almost magnitude-phase form:
  \[
  f_d[n] = \begin{cases} \frac{1}{L}& 0\le m\le L-1\\
    0&\mbox{otherwise}\end{cases}~~~\leftrightarrow~~~
  F_d(\omega)=\left(\frac{\sin(\omega L/2)}{L\sin(\omega/2)}\right)e^{-j\omega\left(\frac{L-1}{2}\right)}
  \]
  By the way, remember we discovered that
  \[
  F_d(\omega)=e^{-j\omega\left(\frac{L-1}{2}\right)}F_c(\omega)
  \]
  Notice anything?
\end{frame}

\begin{frame}
  \frametitle{DTFT of Local Averaging Filters}

  \begin{itemize}
  \item {\bf Centered local average:}
    \begin{align*}
      f_c[n] &= \begin{cases} \frac{1}{L}& -\left(\frac{L-1}{2}\right)\le n\le\left(\frac{L-1}{2}\right)\\
        0&\mbox{otherwise}\end{cases}\\
      F_c(\omega)&=\frac{\sin(\omega L/2)}{L\sin(\omega/2)}
    \end{align*}
  \item {\bf Causal local average:}
    \begin{align*}
      f_d[n] &= \begin{cases} \frac{1}{L}& 0\le n\le L-1\\
        0&\mbox{otherwise}\end{cases}
      F_d(\omega) &=\frac{\sin(\omega L/2)}{L\sin(\omega/2)}e^{-j\omega\left(\frac{L-1}{2}\right)}
    \end{align*}
  \end{itemize}
\end{frame}

%%%%%%%%%%%%%%%%%%%%%%%%%%%%%%%%%%%%%%%%%%%%
\section[Rectangles]{Rectangular Windows}
\setcounter{subsection}{1}

\begin{frame}
  \frametitle{DTFT of Rectangular Window}

  The local summing filter is just a scaled version of the local
  averaging filter. It is commonly called the ``rectangular window,''
  for reasons that we'll explore in future lectures.
  \begin{itemize}
  \item {\bf Centered rectangular window:}
    \begin{align*}
      w_c[n] &= \begin{cases} 1& -\left(\frac{L-1}{2}\right)\le n\le\left(\frac{L-1}{2}\right)\\
        0&\mbox{otherwise}\end{cases}\\
      W_c(\omega)&=\frac{\sin(\omega L/2)}{\sin(\omega/2)}
    \end{align*}
  \item {\bf Causal rectangular window:}
    \begin{align*}
      w_d[n] &= \begin{cases} 1& 0\le n\le L-1\\
        0&\mbox{otherwise}\end{cases}
      W_d(\omega) &=\frac{\sin(\omega L/2)}{\sin(\omega/2)}e^{-j\omega\left(\frac{L-1}{2}\right)}
    \end{align*}
  \end{itemize}
\end{frame}


\begin{frame}
  \frametitle{Dirichlet form}

  Of all four of these signals, the centered rectangular window has
  the simplest DTFT. The textbook calls this signal the ``Dirichlet
  form:''
  \[
  D_L(\omega) = \frac{\sin(\omega L/2)}{\sin(\omega/2)}
  \]
  That is, exactly, the frequency response of a centered rectangular window:
  \[
  d_L[n] = w_c[n] = \begin{cases}
    1 & -\left(\frac{L-1}{2}\right)\le n\le \left(\frac{L-1}{2}\right)\\
    0 & \mbox{otherwise}
  \end{cases}
  \]
\end{frame}

\begin{frame}
  \frametitle{Dirichlet form}

  Here's what it looks like:
  \centerline{\includegraphics[height=2.5in]{exp/dirichletform.png}}
\end{frame}
  
\begin{frame}
  \frametitle{Dirichlet form}

  Since every local averaging filter is based on Dirichlet form, it's
  worth spending some time to understand it better.
  \[
  D_L(\omega) = \frac{\sin(\omega L/2)}{\sin(\omega/2)}
  \]
  \begin{itemize}
  \item It's equal to zero every time $\omega L/2$ is a multiple of $\pi$.  So
    \[
    D_L\left(\frac{2\pi k}{L}\right)  = 0~~\mbox{for all integers}~k~\mbox{except}~k=0
    \]
  \item At $\omega=0$, the value of $\frac{\sin(\omega
    L/2)}{\sin(\omega/2)}$ is undefined, but it's posssible to prove
    that $\lim_{\omega\rightarrow 0}D_L(\omega)=L$.  To make life
    easy, we'll just define it that way:
    \[
    \mbox{DEFINE:}~~D_L(0)=L
    \]
  \end{itemize}
\end{frame}

\begin{frame}
  \frametitle{Dirichlet form}

  Here's what it looks like:
  \centerline{\includegraphics[height=2.5in]{exp/dirichletform.png}}
\end{frame}
  
\begin{frame}
  \frametitle{Local averaging filter}

  Here's what the centered local averaging filter looks like.  Notice that
  it's just $1/L$ times the Dirichlet form:
  \centerline{\includegraphics[height=2.5in]{exp/centeredaveragingfilter.png}}
\end{frame}
  
%%%%%%%%%%%%%%%%%%%%%%%%%%%%%%%%%%%%%%%%%%%%
\section[Summary]{Summary}
\setcounter{subsection}{1}

\begin{frame}
  \frametitle{Summary}

  The DTFT (discrete time Fourier transform) of any signal is
  $X(\omega)$, given by
  \begin{align*}
    X(\omega) &= \sum_{n=-\infty}^\infty x[n]e^{-j\omega n}\\
    x[n] &= \frac{1}{2\pi}\int_{-\pi}^\pi X(\omega)e^{j\omega n}d\omega
  \end{align*}
  Particular useful examples include:
  \begin{align*}
    f[n]=\delta[n] &\leftrightarrow F(\omega)=1\\
    g[n]=\delta[n-n_0] &\leftrightarrow G(\omega)=e^{-j\omega n_0}
  \end{align*}
\end{frame}

\begin{frame}
  \frametitle{Properties of the DTFT}

  Properties worth knowing  include:
  \begin{enumerate}
    \setcounter{enumi}{-1}
  \item Periodicity: $X(\omega+2\pi)=X(\omega)$
  \item Linearity:
    \[z[n]=ax[n]+by[n]\leftrightarrow Z(\omega)=aX(\omega)+bY(\omega)
    \]
  \item Time Shift: $x[n-n_0]\leftrightarrow e^{-j\omega n_0}X(\omega)$
  \item Frequency Shift: $e^{j\omega_0 n}x[n]\leftrightarrow X(\omega-\omega_0)$
  \item Filtering is Convolution:
    \[
    y[n]=h[n]\ast x[n]\leftrightarrow Y(\omega)=H(\omega)X(\omega)
    \]
  \end{enumerate}
\end{frame}

\begin{frame}
  \frametitle{Summary: DTFTs of Rectangular Windows}
  \begin{itemize}
  \item The {\bf Centered Local Averaging Filter} is $1/L$ times
    the Dirichlet form:
    \[
    f_c[n] = \begin{cases}
      \frac{1}{L} & -\left(\frac{L-1}{2}\right)\le n\le \left(\frac{L-1}{2}\right)\\
      0 & \mbox{otherwise}
    \end{cases}~~~\leftrightarrow~~~
    F_c(\omega)=\frac{\sin(\omega L/2)}{L\sin(\omega/2)}
    \]
  \item The {\bf Causal Local Averaging Filter} is $f_c[n]$, delayed by $\frac{L-1}{2}$ samples:
    \[
    f_d[n] = \begin{cases}
      \frac{1}{L} & 0\le n\le L-1\\
      0 & \mbox{otherwise}
    \end{cases}~~~\leftrightarrow~~~
    F_d(\omega)=\frac{\sin(\omega L/2)}{L\sin(\omega/2)}e^{-j\omega\left(\frac{L-1}{2}\right)}
    \]
  \end{itemize}
\end{frame}

\end{document}
